\section{Experiment Protocol and Ablation Definitions}
\label{app:protocol}

\subsection{P3 budget contract}

The per-Study resource ceilings are:
\begin{itemize}
  \item Provider requests: 200 total, shared by Main and Reporter;
  \item mathematical evaluations: 100;
  \item physical launches: 150;
  \item Public Operations: a separate hard limit derived from the historical Full trace and frozen before new runs.
\end{itemize}
No automatic extension or restart is allowed. The 201st Provider request is prohibited. Budget exhaustion without a qualifying candidate is a within-budget failure; exhaustion with an incomplete required evaluation is classified as incomplete. Neither outcome proves global infeasibility.

\subsection{Fair comparison conditions}

New P3 conditions use the same Core commit, runtime, task contract, initial conditions, FamilySpace, parameter bounds, physical evaluator, and resource ceilings. Each condition uses an independent Store and an empty candidate pool. The historical final candidate is not used to initialize a new ablation.

Where possible, a matched-budget Full prefix is reconstructed from the original trace. If this cannot be reconstructed reliably, historical Full is treated as a high-budget reference rather than a matched comparison. The historical Full result used 150 Provider requests, 67 mathematical evaluations, and 98 physical episodes.

\subsection{Nested ablation interpretation}

No-Math disables all optional mathematical research operations, including pairwise interaction. No-Pairwise-Interaction disables only the pairwise four-point \texttt{interact} operation while retaining the other mathematical operations. These conditions are nested, not independent factorial levels. The No-Pairwise-Interaction result therefore measures the marginal value of pairwise interaction conditional on the remaining mathematical research capability.

\subsection{No-Knowledge alignment}

The historical P3 Full knowledge record references 45 records across 17 entry IDs, while the current Core snapshot reportedly loads four general entries. The knowledge assets must be restored or aligned before a strict Full-versus-No-Knowledge comparison. Otherwise, the experiment measures the current four-entry configuration against zero, and the conclusion must be scoped accordingly.

\subsection{Reporter paired snapshots}

Six historical decision snapshots are evaluated in paired Reporter ON/OFF conditions. Both conditions receive the same fixed numerical facts, task contract, and evidence access. Only the Reporter interpretation is changed. The Main decision is scored blind for evidence use, physical reasoning, decision quality, unsupported claims, and context/token cost.


\subsection{Historical study resource references}

The historical studies provide resource references for planning the new hard ceilings:

\begin{table*}[t]
\centering
\caption{Historical resource references.}
\label{tab:historical-cost}
\small
\begin{tabular}{@{}lrrr@{}}
\toprule
Study & Provider requests & Math evaluations/requests & Physical episodes/steps \\
\midrule
M1 & 48 & 7 & 5 episodes \\
P3 & 150 & 67 & 98 episodes \\
PHS-2D & -- & 303 requests & 3,456,000 physics steps \\
\bottomrule
\end{tabular}
\end{table*}

The PHS-2D study also recorded 47 structural-generation attempts. These counts use the original study-specific definitions and should not be combined across studies as if they were one common benchmark cost. In particular, mathematical Worker integration steps and registered physical launches are distinct accounting quantities.

The historical P3 record includes two finite-scale pairwise interaction studies: \texttt{payout\_offsets\_m.tendon\_1} with \texttt{spans.0.radius\_m}, and remote-segment density with cross-section radius. Both used finite perturbations with strict derivative qualification disabled. The historical bounded directional \texttt{optimize} operation did not produce a defined derivative or a comparable improvement and must not be described as successful gradient descent.
